\chapter{Giới thiệu}
\label{chap:intro}

\section{Đặt vấn đề}

Thị trường tiền mã hóa biến động mạnh và chịu tác động đồng thời của nhiều nhóm yếu tố. Dữ liệu giá và các chỉ báo
kỹ thuật phản ánh động lượng ngắn hạn; dữ liệu on-chain như số địa chỉ hoạt động, số giao dịch, tổng giá trị chuyển
phản ánh mức độ sử dụng mạng lưới; tin tức phản ánh tâm lý và sự kiện vĩ mô. Ba nhóm nguồn này khác nhau về tần suất,
độ nhiễu và tầm nhìn thời gian, nên thường đưa ra tín hiệu trái chiều: chỉ báo kỹ thuật có thể báo quá mua trong khi
tin tức tích cực và dữ liệu on-chain đi ngang.

Các mô hình ngôn ngữ lớn (LLM) mở ra khả năng đọc hiểu đồng thời dữ liệu số và văn bản, và nhiều agent giao dịch dựa
trên LLM đã được đề xuất~\cite{yu2023finmem,yu2024fincon,xiao2024tradingagents,li2024cryptotrade}. Tiêu biểu cho thị
trường tiền mã hóa là CryptoTrade~\cite{li2024cryptotrade}: một agent phân tích dữ liệu giá và on-chain, một agent
phân tích tin tức, sau đó một agent giao dịch đọc cả hai bản phân tích và ra quyết định. Trong cách tiếp cận này, khi
các nguồn mâu thuẫn, việc cân nhắc nguồn nào đáng tin hơn diễn ra ngầm trong một lời nhắc duy nhất: hệ thống không đo
mức mâu thuẫn, không giải thích nguyên nhân và không có cơ chế để các quan điểm đối lập phản biện lẫn nhau.

Song song, các nghiên cứu về tranh biện đa tác tử cho thấy để nhiều phiên bản LLM trao đổi và phản biện qua nhiều vòng
giúp cải thiện tính đúng đắn của câu trả lời~\cite{du2024debate,liang2024divergent}. MADAM-RAG~\cite{wang2025madamrag}
đi xa hơn khi giao mỗi tài liệu nguồn cho một agent riêng và cho các agent tranh biện khi các nguồn mâu thuẫn. Tuy
nhiên các kết quả này chủ yếu ở bài toán hỏi đáp. Câu hỏi đặt ra là: \emph{trong bài toán giao dịch, việc tách mỗi
nguồn cho một agent và xử lý mâu thuẫn một cách tường minh có mang lại lợi ích so với việc để một LLM tự gộp các nguồn
hay không?}

\section{Câu hỏi nghiên cứu}

Đề tài trả lời ba câu hỏi, đều trong điều kiện cùng dữ liệu, cùng môi trường giao dịch và cùng mô hình ngôn ngữ:
\begin{description}[leftmargin=1.2cm,style=nextline]
  \item[RQ1] Kiến trúc tách agent theo nguồn kèm phát hiện mâu thuẫn và tranh biện so với CryptoTrade như thế nào về
  lợi suất, Sharpe và tỷ lệ dự đoán đúng hướng?
  \item[RQ2] Mỗi nguồn dữ liệu (thị trường, on-chain, tin tức) đóng góp gì vào kết quả, và đóng góp đó có thay đổi theo
  giai đoạn thị trường không?
  \item[RQ3] Tranh biện có thực sự thay đổi quyết định và cải thiện kết quả không, hay chỉ làm giảm chỉ số mâu thuẫn?
\end{description}

\section{Đóng góp của đề tài}
\begin{enumerate}
  \item \textbf{Hệ thống đa tác tử có xử lý mâu thuẫn tường minh} gồm ba agent chuyên môn, một Conflict Analyzer đo mâu
  thuẫn bằng chỉ số kết hợp phân kỳ KL và phương sai, phân loại mâu thuẫn và phân tích nguyên nhân gốc, một mô-đun
  tranh biện nhiều vòng với cập nhật độ tự tin tất định và kiểm tra nguồn trích dẫn, và một Mediator tổng hợp có trọng
  số (Chương~\ref{chap:method}).
  \item \textbf{Tái hiện baseline đã công bố}: cài đặt lại CryptoTrade, kiểm chứng trùng khớp với mã nguồn gốc của tác
  giả và với Bảng 2 của bài báo đối với các chiến lược luật (Mục~\ref{sec:baseline}).
  \item \textbf{Đánh giá đối chứng có kiểm soát}: hai hệ thống chạy trên cùng dữ liệu, cùng thông tin đầu vào mỗi ngày,
  cùng môi trường giao dịch và cùng một mô hình chạy cục bộ, kèm các thí nghiệm loại bỏ thành phần (ablation) theo
  nguồn dữ liệu và theo tranh biện (Chương~\ref{chap:exp}).
  \item \textbf{Phân tích thực nghiệm}: tranh biện giảm mâu thuẫn nhưng hiếm khi đổi quyết định; giá trị của từng nguồn
  phụ thuộc giai đoạn thị trường; chênh lệch với CryptoTrade đến chủ yếu từ mức vào lệnh chứ không phải từ chất lượng dự
  đoán hướng.
\end{enumerate}

\section{Phạm vi}
Đề tài giới hạn ở Bitcoin, tần suất quyết định theo ngày và dữ liệu công khai của CryptoTrade (01/2023--01/2024). Hệ
thống đưa ra tín hiệu và được đánh giá trong môi trường mô phỏng; không thực hiện giao dịch thật. Mọi thí nghiệm dùng
một mô hình ngôn ngữ mã nguồn mở chạy cục bộ để đảm bảo chi phí bằng không và khả năng tái lập.

\section{Bố cục báo cáo}
Chương~\ref{chap:background} trình bày cơ sở lý thuyết và các nghiên cứu liên quan. Chương~\ref{chap:method} mô tả
kiến trúc và các công thức của hệ thống. Chương~\ref{chap:exp} trình bày thiết lập thực nghiệm, kết quả và phân tích.
Chương~\ref{chap:conclusion} tổng kết và nêu hướng phát triển.
